		
		\chapter{引言}
		
	 [鼠标左键单击选择该段落，输入替换之。内容为小四号宋体。] 引言（或绪论）简要说明研究工作的目的、范围、相关领域的前人工作和知识空白、理论基础和分析、研究设想、研究方法和实验设计、预期结果和意义等。应言简意赅，不要与摘要雷同，不要成为摘要的注释。一般教科书中有的知识，在引言中不必赘述。
	 
	 \section{参考文献示例}
	 参考文献bib文件放在reference文件夹中，在主文件demo中引用了bib文件，在所有chap中都可以直接cite。比如\cite{silver2016mastering}
	 
	 \section{算法框示例}
	 
	 \begin{algorithm}[h]
	 	\caption{ 优化算法}
	 	\label{gtdmc:alg:minmax}			
	 	\KwIn{ 迭代次数$ T$，   待学习参数初始值$x_1$ 、 $y_1$ ， 学习率 $\alpha$ }
	 	\For{$ t = 1, \dots, T $}{
	 		更新参数:  
	 		
	 		$ y_{t+1} = \mathcal{P}_{\mathcal{X}_y}\left(y_t + \alpha_t(\hat{b}_t - \hat{A}_t x_t -\hat{M}_ty_t)\right)  $
	 		
	 		$ x_{t+1} = \mathcal{P}_{\mathcal{X}_x}\left(x_t + \alpha_t\hat{A}_t^\top y_t\right) $
	 	}
	 	
	 	
	 	\KwOut{ 		$ \quad 	 \tilde{x}_T = \frac{\sum_{t=1}^{T}\alpha_t x_t}{\sum_{t=1}^{T}\alpha_t}  \qquad \tilde{y}_T = \frac{\sum_{t=1}^{T}\alpha_t y_t}{\sum_{t=1}^{T}\alpha_t} $	 }				
	 	
	 	
	 \end{algorithm}	 
 
 \section{图片示例}
 
 
 
 			\begin{figure}[h]
 	\centering
 	\subfloat[a]{
 		\label{figa}
 		\includegraphics[width=1.8in,height=1.5in]{a.png}}
 	\subfloat[b]{
 		\label{figc}
 		\includegraphics[width=1.8in,height=1.5in]{a.png}} 
 	\subfloat[c ]{
 		\label{fige}
 		\includegraphics[width=1.8in,height=1.5in]{a.png}}	
 	
 	
 	\subfloat[d ]{
 		\label{figb}
 		\includegraphics[width=1.8in,height=1.5in]{a.png}}
 	\label{Fig1}	
 	\subfloat[e ]{
 		\label{figd}
 		\includegraphics[width=1.8in,height=1.5in]{a.png}}	
 	\subfloat[f ]{
 		\label{figf}
 		\includegraphics[width=1.8in,height=1.5in]{a.png}}		
 	\bicaption{你好，世界}{Hello world } 
 	\label{fig1}
 \end{figure}	
 